% (c) Nikita Lisitsa, lisyarus@gmail.com, 2026

\documentclass[10pt]{beamer}

\usepackage[T2A]{fontenc}
\usepackage[russian]{babel}
\usepackage{minted}

\usepackage[most]{tcolorbox}
\tcbuselibrary{minted}

\usepackage{graphicx}
\graphicspath{ {./images/} }

\usepackage{adjustbox}

\usepackage{color}
\usepackage{soul}

\usepackage{hyperref}

\usepackage{amsmath}

\usepackage{tikz}
\usetikzlibrary{decorations}
\usetikzlibrary{decorations.pathreplacing}
\usepackage{xifthen}

\usetheme{metropolis}
\setminted{fontsize=\footnotesize}

\definecolor{red}{rgb}{1,0,0}
\definecolor{green}{rgb}{0,0.5,0}
\definecolor{blue}{rgb}{0,0,1}
\definecolor{magenta}{rgb}{0.75,0,0.75}
\definecolor{codebg}{RGB}{29,35,49}

\makeatletter
\newcommand{\slideimage}[1]{
  \begin{figure}
    \begin{adjustbox}{width=\textwidth, totalheight=\textheight-2\baselineskip-2\baselineskip,keepaspectratio}
      \includegraphics{#1}
    \end{adjustbox}
  \end{figure}
}
\makeatother

\usemintedstyle{tango}

\title{Компьютерная графика}
\subtitle{Лекция 8: Геометрические шейдеры, compute шейдеры, storage buffers, shader atomics, shadow volumes, shadow mapping}
\date{2026}

\setbeamertemplate{footline}[frame number]

\begin{document}

\usemintedstyle{solarized-light}

\frame{\titlepage}

\begin{frame}[fragile]
\frametitle{Виды шейдеров}
\begin{itemize}
\item \textbf{Render pipeline}
\begin{itemize}
\item Vertex shader
\item Tessellation control shader
\item Tessellation evaluation shader
\item Geometry shader
\item Fragment shader
\end{itemize}
\pause
\item \textbf{Mesh shading pipeline}
\begin{itemize}
\item Task shader
\item Mesh shader
\item Fragment shader
\end{itemize}
\pause
\item \textbf{Ray tracing pipeline}
\begin{itemize}
\item Ray generation shader
\item Intersection shader
\item Any-hit shader
\item Closest-hit shader
\item Miss shader
\end{itemize}
\pause
\item \textbf{Compute pipeline}
\begin{itemize}
\item Compute shader
\end{itemize}
\end{itemize}
\end{frame}

\begin{frame}[fragile]
\frametitle{Виды шейдеров}
\begin{itemize}
\item Из-за кросс-платформенности, WebGPU поддерживает только базовый \textbf{render pipeline (vertex + fragment)} и \textbf{compute pipeline}
\pause
\item В wgpu есть экспериментальная поддержка \textbf{mesh shaders} и \textbf{ray queries}, ведётся работа над поддержкой \textbf{ray tracing pipelines}
\end{itemize}
\end{frame}

\begin{frame}[fragile]
\frametitle{Геометрические шейдеры}
\begin{itemize}
\item Опциональный шейдер, встраивающийся после вершинного шейдера, до perspective divide
\pause
\item Оперирует целыми примитивами (точками/линиями/треугольниками), т.е. наборами вершин
\pause
\item Может на лету менять тип примитива и количество вершин
\pause
\item Поддерживаются в OpenGL, DirectX, Vulkan
\pause
\item \textbf{Не поддерживается} в WebGPU, Metal
\end{itemize}
\end{frame}

\begin{frame}[fragile]
\frametitle{Геометрические шейдеры}
\begin{itemize}
\item Геометрические шейдеры тяжело ложатся на архитектуру GPU из-за отсутствия чёткого соответствия входов и выходов
\pause
\item Требуют дополнительной временной памяти, не используют вершинный кеш
\pause
\item Производительность очень сильно зависит от GPU, драйвера, и самого шейдера
\pause
\item В современности их стараются не использовать, заменяя более надёжными инструментами: более сложными вершинными шейдерами, instancing'ом, compute шейдерами
\end{itemize}
\end{frame}

\begin{frame}<1>[fragile,label=geometry_shader_examples]
\frametitle{Геометрические шейдеры: примеры использования}
\fontsize{10pt}{10pt}
\begin{itemize}
\item Расчёт нормалей для flat shading'а
\pause
\begin{itemize}
\item Вход: треугольник (тройка вершин)
\item Выход: те же вершины с посчитанной нормалью
\end{itemize}
\pause
\item Визуализация нормалей
\pause
\begin{itemize}
\item Вход: вершина с нормалью
\item Выход: линия из двух вершин -- исходная вершина, исходная вершина + нормаль
\end{itemize}
\pause
\item Shadow volumes -- алгоритм рисования теней
\pause
\item Billboards -- плоские фигуры, всегда смотрящие в сторону камеры
\pause
\begin{itemize}
\item Вход: одна вершина (точка)
\item Выход: набор треугольников
\pause
\item Системы частиц
\pause
\item Облака
\pause
\item Далёкие деревья
\pause
\item Трава
\end{itemize}
\end{itemize}
\end{frame}

\begin{frame}[fragile]
\frametitle{Flat shading}
\slideimage{flat-shading.png}
\end{frame}

\againframe<1-3>{geometry_shader_examples}

\begin{frame}[fragile]
\frametitle{Визуализация нормалей}
\slideimage{normals-viz.png}
\end{frame}

\againframe<3-6>{geometry_shader_examples}

\begin{frame}[fragile]
\frametitle{Billboards}
\slideimage{billboards.jpg}
\end{frame}

\againframe<6-8>{geometry_shader_examples}

\begin{frame}[fragile]
\frametitle{Billboards: система частиц (дым)}
\slideimage{smoke.jpg}
\end{frame}

\againframe<8-9>{geometry_shader_examples}

\begin{frame}[fragile]
\frametitle{Billboards: облака}
\slideimage{clouds.jpg}
\end{frame}

\againframe<9-10>{geometry_shader_examples}

\begin{frame}[fragile]
\frametitle{Billboards: деревья}
\slideimage{trees.jpg}
\end{frame}

\againframe<10-11>{geometry_shader_examples}

\begin{frame}[fragile]
\frametitle{Billboards: трава}
\slideimage{grass.jpg}
\end{frame}

\begin{frame}[fragile]
\frametitle{Геометрические шейдеры: ссылки}
\begin{itemize}
\item \href{https://www.khronos.org/opengl/wiki/Geometry_Shader}{\nolinkurl{khronos.org/opengl/wiki/Geometry\_Shader}}
\item \href{https://learnopengl.com/Advanced-OpenGL/Geometry-Shader}{\nolinkurl{learnopengl.com/Advanced-OpenGL/Geometry-Shader}}
\item \href{https://open.gl/geometry}{\nolinkurl{open.gl/geometry}}
\item \href{https://www.lighthouse3d.com/tutorials/glsl-tutorial/geometry-shader}{\nolinkurl{lighthouse3d.com/tutorials/glsl-tutorial/geometry-shader}}
\item \href{https://developer.nvidia.com/gpugems/gpugems/part-i-natural-effects/chapter-7-rendering-countless-blades-waving-grass}{\texttt{GPU Gems, Chapter 7. Rendering Countless Blades of Waving Grass}}
\end{itemize}
\end{frame}

\begin{frame}[fragile]
\frametitle{Compute шейдеры}
\begin{itemize}
\item Compute шейдеры -- механизм для запуска произвольных параллельных вычислений на GPU, минующий весь графический конвейер
\pause
\item Для их запуска нужен не render pass, а \textit{compute pass} (создаётся так же через command encoder)
\pause
\item Он не привязан к цветовым буферам или буферам глубины, для его создания не нужно почти никаких настроек
\pause
\item Так же как с render pass'ом, можно выставить bind группы и immediates
\end{itemize}
\end{frame}

\begin{frame}[fragile]
\frametitle{Compute pipeline}
\begin{itemize}
\item Compute pass умеет использовать только \textit{compute pipelines}
\pause
\item Compute pipeline создаётся аналогично render pipeline'ам: нужен pipeline layout, shader module и entry point
\pause
\item Так как compute шейдеры не используют графический конвейер, связанных с ним настроек нет (vertex/fragment/primitive state и т.п.)
\pause
\item Команд \mintinline{cpp}|draw()| у compute pass'а нет, основной способ инициировать вычисление -- \mintinline{cpp}|compute_pass.dispatch_workgroups()|
\end{itemize}
\end{frame}

\begin{frame}[fragile]
\frametitle{Compute шейдеры: модель исполнения}
\begin{itemize}
\item Compute shader'ы запускаются сеткой \(N_x\times N_y\times N_z\) ячеек, каждая ячейка соответствует одному потоку исполнения внутри GPU
\pause
\item GPU исполняют задачи SIMD-блоками (обычно 32 или 64 парралельных исполнения)
\pause
\item Compute шейдеры явно отражают это в своём интерфейсе: они запускаются как набор из нескольких \textit{workgroup} одинакового размера \(S_x\times S_y\times S_z\) (\mintinline{cpp}|workgroup_size|, задаётся в шейдере)
\pause
\item Аргументы к \mintinline{cpp}|compute_pass.dispatch_workgroups()| -- количество запущенных workgroup \(C_x\times C_y\times C_z\) (\mintinline{cpp}|workgroup_count|)
\pause
\item В итоге, \(N_x = S_x \cdot C_x\), и аналогично по другим осям
\end{itemize}
\end{frame}

\begin{frame}[fragile]
\frametitle{Compute шейдеры: модель исполнения}
\begin{itemize}
\item Выбор размера workgroup -- сложная задача:
\pause
\begin{itemize}
\item Один SIMD-блок GPU может выполнять только одну workgroup \(\Longrightarrow\) в идеале размер workgroup должен быть кратен размеру SIMD-блока
\pause
\item Большой размер workgroup требует больше локальной памяти, но может позволить переиспользовать больше вычислений (через workgroup shared memory)
\pause
\item Как слишком маленькими, так и слишком большими workgroup можно упереться в различные hardware лимиты
\end{itemize}
\pause
\item Почти всегда SIMD-блоки имеют размер в 32 или 64 потока, поэтому стандартный выбор для workgroup size -- 64 (произведение по всем осям)
\pause
\item Это удобно для обработки изображений (блоки \(8\times 8\)) или 3D-данных (блоки \(4\times 4\times 4\))
\pause
\item Если реальный объём работы (например, размер изображения) не кратен размеру workgroup, лишние вычисления обычно явно пропускаются в шейдере
\end{itemize}
\end{frame}

\begin{frame}[fragile]
\frametitle{Compute шейдеры: детали}
\begin{itemize}
\item Entry point для compute pipeline (функция в шейдере) должна быть помечена атрибутами \mintinline{wgsl}|@compute| и \mintinline{wgsl}|@workgroup_size(x, y, z)| (\mintinline{wgsl}|y| и \mintinline{wgsl}|z| можно не указывать, если они равны 1)
\pause
\item Встроенные переменные compute шейдера:
\begin{itemize}
\item \mintinline{wgsl}|@builtin(local_invocation_id)| -- 3D индекс потока внутри workgroup (\mintinline{cpp}|0..workgroup_size-1|)
\item \mintinline{wgsl}|@builtin(workgroup_id)| -- 3D индекс workgroup внутри вызова dispatch (\mintinline{cpp}|0..workgroup_count-1|)
\item \mintinline{wgsl}|@builtin(global_invocation_id)| -- 3D индекс потока внутри вызова dispatch (вычисляется как \mintinline{cpp}|workgroup_id * workgroup_size + local_invocation_id|)
\item \mintinline{wgsl}|@builtin(local_invocation_index)| -- линеаризованный индекс потока внутри workgroup
\item \mintinline{wgsl}|@builtin(global_invocation_index)| -- линеаризованный индекс потока внутри dispatch
\end{itemize}
\pause
\item Встроенных входных или выходных данных нет, все операции делаются через данные, переданные в bind groups
\end{itemize}
\end{frame}

\begin{frame}[fragile]
\frametitle{Compute шейдеры: пример}
\usemintedstyle{lightbulb}
\begin{tcblisting}{listing engine=minted, minted language=wgsl, minted options={fontsize=\scriptsize}, listing only, colback=codebg, colframe=codebg, rounded corners}
@group(0) @binding(0) var<storage, read> a: array<f32>;
@group(0) @binding(1) var<storage, read> b: array<f32>;
@group(0) @binding(2) var<storage, read_write> c: array<f32>;

@compute @workgroup_size(64)
fn main(@builtin(global_invocation_id) id: vec3u) {
    let index = id.x;
    if (index < arrayLength(&c)) {
        c[index] = a[index] + b[index];
    }
}
\end{tcblisting}
\end{frame}

\begin{frame}[fragile]
\frametitle{Compute шейдеры: пример}
\usemintedstyle{lightbulb}
\begin{tcblisting}{listing engine=minted, minted language=wgsl, minted options={fontsize=\scriptsize}, listing only, colback=codebg, colframe=codebg, rounded corners}
@group(0) @binding(0) var image:
  texture_storage_2d<rgba8unorm, read_write>;

@compute @workgroup_size(8, 8, 1)
fn main(@builtin(global_invocation_id) id: vec3u) {
    if (any(id.xy >= textureDimensions(image))) {
      return;
    }

    let colorIn = textureLoad(image, id.xy);

    let colorOut = vec4f(vec3f(1.0) - colorIn.rgb, colorIn.a);

    textureStore(image, id.xy, colorOut);
}
\end{tcblisting}
\end{frame}

\begin{frame}[fragile]
\frametitle{Storage buffers}
\begin{itemize}
\item Буферы, созданные с флагом \mintinline{cpp}|usage = STORAGE|
\pause
\item Их можно добавить в bind группу с режимом \mintinline{cpp}|storage| или \mintinline{cpp}|read_only_storage|
\pause
\item В шейдере должны быть объявлены как \mintinline{wgsl}|var<storage, read>| или \mintinline{wgsl}|var<storage, read_write>|, соответственно
\pause
\item Тип данных -- любой, в том числе может быть массив неизвестного размера: \mintinline{wgsl}|array<vec4f>| (размер массива можно узнать через функцию \mintinline{wgsl}|arrayLength|)
\pause
\item Основной способ передать произвольные данные на GPU, размер ограничен размером видеопамяти
\pause
\item Не read-only буферы можно использовать только во фрагментных и compute шейдерах (с фичей \mintinline{cpp}|VERTEX_WRITABLE_STORAGE| можно и в вершинных -- не поддерживается в браузерной версии WebGPU)
\end{itemize}
\end{frame}

\begin{frame}[fragile]
\frametitle{Storage textures}
\begin{itemize}
\item Текстуры, созданные с флагом \mintinline{cpp}|usage = STORAGE_BINDING|
\pause
\item Их тоже можно добавить в bind группу с явным указанием формата
\pause
\item В шейдере тоже объявляются с явным указанием формата, например \mintinline{wgsl}|texture_storage_2d<rgba8unorm, read_write>|
\pause
\item Нельзя использовать вместе с семплерами, только читать и писать конкретные пиксели через \mintinline{wgsl}|textureLoad| и \mintinline{wgsl}|textureStore|
\pause
\item Поддерживают ограниченный набор форматов
\end{itemize}
\end{frame}

\begin{frame}[fragile]
\frametitle{Тени: теория}
\begin{itemize}
\item Точка сцены, в которую не попадает (заблокирован чем-то) прямой свет из конкретного источника света
\item Свойство точки по отношению к конкретному источнику света
\end{itemize}
\end{frame}

\begin{frame}[fragile]
\frametitle{Тени}
\slideimage{shadows1.png}
\end{frame}

\begin{frame}[fragile]
\frametitle{Тени: теория}
\begin{itemize}
\item Если источник света точечный (или бесконечно удалённый), тень -- бинарное свойство: луч из точки сцены в источник света или пересекает что-то (точка в тени), или нет (точка не в тени) \begin{math}\Longrightarrow\end{math} жёсткие тени (\textit{hard shadows})
\pause
\item Если источник света объёмный, точка сцены может находиться в тени относительно части источника света, и не находиться в тени относительно другой его части \begin{math}\Longrightarrow\end{math} мягкие тени (\textit{soft shadows})
\pause
\begin{itemize}
\item Точки, полностью находящиеся в тени -- \textit{umbra}
\item Точки, частично находящиеся в тени -- \textit{penumbra} (\textit{полутень})
\end{itemize}
\end{itemize}
\end{frame}

\begin{frame}[fragile]
\frametitle{Тени: мягкие vs жёсткие}
\slideimage{shadow-scheme1.png}
\end{frame}

\begin{frame}[fragile]
\frametitle{Мягкие тени}
\slideimage{shadows2.png}
\end{frame}

\begin{frame}[fragile]
\frametitle{Солнечное затмение}
\slideimage{eclipse.png}
\end{frame}

\begin{frame}[fragile]
\frametitle{Тени: теория}
\begin{itemize}
\item Реальные источники света -- объёмные
\pause
\item Размер и форма penumbra зависит от размеров и формы источника света, а также от расстояния до него (чем дальше, тем меньше penumbra)
\pause
\begin{itemize}
\item В пределе расстояния \begin{math}\rightarrow\infty\end{math} мягкая тень вырождается в жёсткую
\end{itemize}
\pause
\item В real-time графике получить правильные тени (как жёсткие, так и мягкие) довольно сложно
\end{itemize}
\end{frame}

\begin{frame}[fragile]
\frametitle{Тени: raytracing}
\begin{itemize}
\item При простом raytracing'е (\textit{Whitted-style}, без полного решения rendering equation) для получения жёсткой тени можно послать дополнительный луч в направлении света (shadow ray): если он что-то пересёк, точка находится в тени
\pause
\item При raytracing'е с монте-карло интегрированием для global illumination тени (как жёсткие, так и мягкие) получаются автоматически
\end{itemize}
\end{frame}

\begin{frame}[fragile]
\frametitle{Тени: real-time алгоритмы}
\begin{itemize}
\item Shadow volumes
\item Shadow mapping
\item Raytraced shadows
\end{itemize}
\end{frame}

\begin{frame}<1>[fragile,label=shadow_volumes]
\frametitle{Shadow volumes (stencil shadows)}
\begin{itemize}
\item Тень от конкретного объекта -- некая (полубесконечная) трёхмерная область пространства (shadow volume)
\pause
\item Нарисуем её как трёхмерный объект (построим и нарисуем её грани)
\pause
\item Если в конкретный пиксель попало одинаковое количество front-facing и back-facing граней, то он не в тени, иначе -- в тени
\pause
\item Количество пересечений луча из камеры с гранями shadow volume вычисляем с помощью stencil буфера
\end{itemize}
\end{frame}

\begin{frame}<1>[fragile]
\frametitle{Shadow volumes (stencil shadows)}
\slideimage{shadow-volume.png}
\end{frame}

\againframe<1->{shadow_volumes}

\begin{frame}[fragile]
\frametitle{Shadow volumes (stencil shadows)}
\slideimage{shadow-volumes-scheme.png}
\end{frame}

\begin{frame}[fragile]
\frametitle{Shadow volumes: алгоритм}
\begin{itemize}
\item Выключаем рисование в цветовой буфер и буфер глубины (как настройки render pipeline, или как отдельный render pass), но \textbf{не выключаем} тест глубины
\pause
\item Настраиваем запись в stencil буфер: \mintinline{cpp}|+1| для front-facing треугольников, \mintinline{cpp}|-1| для back-facing
\pause
\item Вычисляем и рисуем shadow volume для каждого объекта
\pause
\item Обратно включаем рисование цвета и глубины
\pause
\item Настраиваем stencil-тест так, чтобы рисовать только пиксели, в которых \mintinline{cpp}|stencil == 0|, и рисуем сцену шейдером, вычисляющим освещение
\pause
\item Настраиваем stencil-тест так, чтобы рисовать только пиксели, в которых \mintinline{cpp}|stencil != 0|, и рисуем сцену шейдером, игнорирующим освещение
\pause
\item Повторяем для каждого источника света
\end{itemize}
\end{frame}

\begin{frame}[fragile]
\frametitle{Shadow volume}
\slideimage{shadow-volume2.jpg}
\end{frame}

\begin{frame}[fragile]
\frametitle{Shadow volume}
\slideimage{shadow-volume3.png}
\end{frame}

\begin{frame}[fragile]
\frametitle{Shadow volumes (Doom 3)}
\slideimage{doom3-shadows.png}
\end{frame}

\begin{frame}[fragile]
\frametitle{Shadow volumes: как вычислять shadow volume?}
\begin{itemize}
\item Шаг 1: определить shadow edges -- рёбра модели, у которых один соседний треугольник обращён к источнику света \begin{math}\left(\vec n \cdot \vec r > 0\right)\end{math}, а второй -- против источника света \begin{math}\left(\vec n \cdot \vec r \leq 0\right)\end{math}
\pause
\item Щаг 2: для каждого shadow edge построить четырёхугольную грань shadow volume -- к точкам ребра \begin{math}p_1, p_2\end{math} добавлются точки \begin{math}p_1 + D \cdot \frac{p_1 - o}{|p_1 - o|}\end{math} и аналогично для \begin{math}p_2\end{math}
\begin{itemize}
\item \begin{math}o\end{math} -- координаты источника света
\item \begin{math}D\end{math} -- расстояние, до которого отбрасывается тень (в идеале мы хотим \begin{math}D\rightarrow\infty\end{math}, для этого можно воспользоваться однородными координатами)
\end{itemize}
\end{itemize}
\end{frame}

\begin{frame}[fragile]
\frametitle{Shadow volumes (Doom 3)}
\slideimage{shadow-edge.jpg}
\end{frame}

\begin{frame}[fragile]
\frametitle{Shadow volumes: как вычислять shadow volume?}
\begin{itemize}
\item Два варианта: на CPU или на GPU
\pause
\item На CPU: очень дорого (положение модели и источника света могут меняться каждый кадр)
\pause
\item На GPU -- много вариантов (вершинный/геометрический/compute шейдер), главная проблема -- передать в шейдер две вершины одного ребра + две вершины соседних треугольников
\pause
\item В любом случае нужны входные данные, специально адаптированные для генерации shadow volume
\pause
\item Шейдер проверяет, является ли ребро shadow edge, и генерирует грань shadow volume
\end{itemize}
\end{frame}

\begin{frame}[fragile]
\frametitle{Shadow volumes}
\begin{itemize}
\item {\color{green}+} Pixel-perfect жёсткие тени
\item {\color{red}---} Алиасинг
\item {\color{red}---} Сложно получить мягкие тени
\item {\color{red}---} Нужно на лету генерировать геометрию
\item {\color{red}---} Требует правильной входной геометрии (без дырок и дублирования)
\item {\color{red}---} Не работает, если камера находится внутри тени (исправляется с помощью т.н. Carmack's Reverse)
\item {\color{red}---} Ресурсозатратен: даже для маленького объекта его shadow volume может занимать значительную часть сцены \begin{math}\Longrightarrow\end{math} слишком много растеризации, чтения/записи памяти (stencil buffer)
\end{itemize}
\end{frame}

\begin{frame}[fragile]
\frametitle{Shadow volumes с плохой геометрией}
\slideimage{shadow-volume-bad-mesh.jpg}
\end{frame}

\begin{frame}[fragile]
\frametitle{Shadow volumes}
\begin{itemize}
\item Thief 3
\item Doom 3
\item Quake 4
\item Prey
\item Far Cry
\item F.E.A.R 1, 2, 3
\item S.T.A.L.K.E.R.
\item Timeshift
\item и др.
\end{itemize}
\end{frame}

\begin{frame}[fragile]
\frametitle{Shadow volumes: ссылки}
\begin{itemize}
\item \href{https://en.wikipedia.org/wiki/Shadow_volume}{\nolinkurl{en.wikipedia.org/wiki/Shadow\_volume}}
\item \href{https://ogldev.org/www/tutorial40/tutorial40.html}{\nolinkurl{ogldev.org/www/tutorial40/tutorial40.html}}
\item \href{https://developer.nvidia.com/gpugems/gpugems/part-ii-lighting-and-shadows/chapter-9-efficient-shadow-volume-rendering}{\texttt{GPU Gems, Chapter 9. Efficient Shadow Volume Rendering}}
\end{itemize}
\end{frame}

\begin{frame}[fragile]
\frametitle{Shadow mapping: идея}
\begin{itemize}
\item Самое популярное семейство алгоритмов рисования теней
\pause
\item Идея: представим, что источник света -- это камера (наблюдатель)
\item Тень -- это то, что \textit{не видит} источник света
\pause
\item Как понять, что видит источник света? \pause Нарисовать!
\end{itemize}
\end{frame}

\begin{frame}[fragile]
\frametitle{Shadow mapping: алгоритм}
\begin{itemize}
\item Рисуем сцену в текстуру (\textit{shadow map}, карта теней) с точки зрения источника света (цветовой буфер не нужен, достаточно буфера глубины)
\pause
\item Рисуем сцену на экран, фрагментный шейдер читает \textit{буфер глубины} (shadow map), нарисованный с точки зрения источника света -- там содержатся расстояния от источника до ближайшей (к источнику) поверхности
\item Если расстояние от источника до рисуемого пикселя больше, чем значение из буфера глубины, -- текущий пиксель находится в тени
\pause
\item \textbf{\alert{N.B.}} Линейная фильтрация для буфера глубины \underline{не имеет смысла} -- будут артефакты на границах объектов
\end{itemize}
\end{frame}

\begin{frame}[fragile]
\frametitle{Shadow mapping: алгоритм}
\slideimage{shadow-mapping1.png}
\end{frame}

\begin{frame}[fragile]
\frametitle{Shadow mapping: фрагментный шейдер}
\usemintedstyle{lightbulb}
\begin{tcblisting}{listing engine=minted, minted language=wgsl, minted options={fontsize=\scriptsize}, listing only, colback=codebg, colframe=codebg, rounded corners}
@group(0) @binding(0) var<uniform> shadow_projection: mat4x4f;
@group(0) @binding(1) var shadow_map: texture_depth_2d;
@group(0) @binding(2) var shadow_map_sampler: sampler;

@fragment
fn fragmentMain(in: VertexOut) -> @location(0) vec4f {
  var shadow_coord = shadow_projection * in.world_position;
  shadow_coord /= shadow_coord.w; // perspective divide

  var is_in_shadow = false;
  if (all(abs(shadow_coord.xy) < vec2f(1.0))) {
    let shadow_texcoord = vec2f(0.5)
      + shadow_coord.xy * vec2f(0.5, -0.5);

    is_in_shadow = textureSample(shadow_map, shadow_map_sampler,
      shadow_texcoord) < shadow_coord.z;
  }
}
\end{tcblisting}
\end{frame}

\begin{frame}[fragile]
\frametitle{Shadow mapping: буфер глубины}
\slideimage{shadow-mapping-depth.png}
\end{frame}

\begin{frame}[fragile]
\frametitle{Shadow mapping: результат}
\slideimage{shadow-mapping2.png}
\end{frame}

\begin{frame}[fragile]
\frametitle{Shadow mapping: проекция}
\begin{itemize}
\item Как выбрать проекцию камеры (\mintinline{cpp}|shadow_projection|) для shadow map?
\pause
\item Удалённый источник освещает всю сцену с одного направления \begin{math}\Longrightarrow\end{math} удобно использовать ортографическую проекцию
\pause
\item Проекцию надо подобрать так, чтобы в видимую область попала вся сцена
\pause
\item Можно спроецировать все объекты на плоскость, перпендикулярную направлению на свет (которое можно считать Z-осью этой камеры), взять любые два ортогональных вектора X и Y в этой плоскости, и найти bounding box сцены в координатах XYZ -- это даст размеры и ориентацию области shadow map
\end{itemize}
\end{frame}

\begin{frame}[fragile]
\frametitle{Shadow mapping: проекция}
\begin{itemize}
\item Точечный источник светит во все стороны \begin{math}\Longrightarrow\end{math} удобно использовать перспективную проекцию и cubemap-текстуру
\pause
\item 6 shadow map (по одной на каждую грань) с соответствующей перспективной проекцией
\pause
\item Положение камеры совпадает с положением источника света, виртуальный `экран' совпадает с отрисовываемой гранью shadow map
\end{itemize}
\end{frame}

\begin{frame}[fragile]
\frametitle{Shadow acne}
\slideimage{shadow-acne.png}
\end{frame}

\begin{frame}<1>[fragile,label=acne]
\frametitle{Shadow acne}
\begin{itemize}
\item Артефакт, возникающий из-за дискретности shadow map: если луч света не перпендикулярен плоскости объекта, рядом с центром пикселя shadow map будут точки с меньшей или большей глубиной
\pause
\item Первый способ исправить: добавить к значению, прочитанному из из shadow map, некую константу (shadow bias)
\begin{itemize}
\item Часто эту константу умножают на тангенс угла падения света
\end{itemize}
\pause
\item Второй способ исправить: рисовать только back-facing грани (т.е. обращённые против света) в shadow map
\begin{itemize}
\item Теперь shadow acne возникнет на back-facing гранях, но мы и без shadow map знаем, что они в тени!
\item Работает только с правильной геометрией (двусторонней, без дырок, и т.п.)
\end{itemize}
\end{itemize}
\end{frame}

\begin{frame}[fragile]
\frametitle{Shadow acne}
\slideimage{shadow-acne-scheme.png}
\end{frame}

\againframe<1->{acne}

\begin{frame}[fragile]
\frametitle{Back-facing shadows}
\slideimage{shadow-mapping-backface.png}
\end{frame}

\begin{frame}[fragile]
\frametitle{Peter panning}
\slideimage{peter-panning.png}
\end{frame}

\begin{frame}[fragile]
\frametitle{Peter panning}
\begin{itemize}
\item Объекты кажутся приподнятыми над поверхностью из-за того, что их тень начинается не сразу под объектом
\pause
\item Усиливается при использовании shadow bias
\pause
\item Лучший способ исправить -- не использовать слишком тонкую геометрию
\end{itemize}
\end{frame}

\begin{frame}[fragile]
\frametitle{Shadow mapping}
\slideimage{shadow-mapping3.png}
\end{frame}

\begin{frame}[fragile]
\frametitle{Разрешение shadow map}
\slideimage{shadow-map-resolution.png}
\end{frame}

\begin{frame}[fragile]
\frametitle{Shadow mapping}
\fontsize{10pt}{10pt}
\begin{itemize}
\item При низком разрешении shadow map хорошо видны `пиксели' теней
\pause
\item Чем выше разрешение, тем лучше жёсткая тень
\pause
\begin{itemize}
\item Алгоритм silhouette mapping позволяет получить жёсткую тень без слишком большого разрешения
\end{itemize}
\pause
\item Можно слегка размыть тень, избавившись от `пикселей' и получив аппроксимацию мягких теней
\pause
\begin{itemize}
\item Размывать нужно результат сравнения глубины со значением в shadow map, а \textbf{не} саму shadow map!
\pause
\item Этот приём известен как \textit{percentage-closer filtering (PCF)}
\pause
\item В WebGPU есть встроенная реализация в виде \textit{comparison sampling}
\end{itemize}
\pause
\item Более современные вариации shadow maps пишут в текстуру не глубину, а какие-нибудь значения, которые имеет смысл складывать и усреднять
\pause
\begin{itemize}
\item Позволяет использовать фильтрацию текстур, mipmaps, размытие shadow map
\pause
\item Exponential shadow maps, variance shadow maps
\end{itemize}
\end{itemize}
\end{frame}

\begin{frame}[fragile]
\frametitle{Comparison sampling}
\begin{itemize}
\item Нужно создать специальный sampler с настройкой \mintinline{cpp}|compare| -- она будет управлять тем, какой тип сравнения будет происходить (\mintinline{cpp}|Less|, \mintinline{cpp}|Equal|, и т.п.)
\pause
\item В bind group layout нужно указать \mintinline{cpp}|sampler_binding_type = Comparison|
\pause
\item В шейдере нужно объявить uniform-переменную типа \mintinline{wgsl}|sampler_comparison|
\pause
\item В шейдере используем функцию \mintinline{wgsl}|textureSampleCompare()| -- она сравнивает значения глубин соседних пикселей с переданной глубиной, и возвращает результат билинейной интерполяции результатов сравнения
\end{itemize}
\end{frame}

\begin{frame}[fragile]
\frametitle{Shadow sampler с линейной фильтрацией}
\slideimage{pcf.png}
\end{frame}

\begin{frame}[fragile]
\frametitle{Shadow sampler с линейной фильтрацией и 7x7 размытием по Гауссу}
\slideimage{pcf_gauss.png}
\end{frame}

\begin{frame}[fragile]
\frametitle{Shadow mapping: ссылки}
\begin{itemize}
\item \href{https://learnopengl.com/Advanced-Lighting/Shadows/Shadow-Mapping}{\nolinkurl{learnopengl.com/Advanced-Lighting/Shadows/Shadow-Mapping}}
\item \href{http://www.opengl-tutorial.org/intermediate-tutorials/tutorial-16-shadow-mapping}{\nolinkurl{opengl-tutorial.org/intermediate-tutorials/tutorial-16-shadow-mapping}}
\item \href{https://ogldev.org/www/tutorial23/tutorial23.html}{\nolinkurl{ogldev.org/www/tutorial23/tutorial23.html}}
\item \href{https://developer.nvidia.com/gpugems/gpugems2/part-ii-shading-lighting-and-shadows/chapter-17-efficient-soft-edged-shadows-using}{\texttt{GPU Gems 2, Chapter 17. Efficient Soft-Edged Shadows Using Pixel Shader Branching}}
\end{itemize}
\end{frame}

\end{document}
